\chapter{Gradient Computation via AD in Chaotic Systems}
\label{chap:lorenz63}

To understand the challenges of calibrating a climate model against reference data, it is useful to begin with a low-dimensional system that exhibits a key feature of atmospheric dynamics: deterministic chaos. We first introduce the mathematical background of chaotic dynamical systems (Section~\ref{sec:chaos}), then present the Lorenz-63 system as a model for such behavior (Section~\ref{sec:lorenz63}). We continue to outline the automatic differentiation methodology used throughout this thesis (Section~\ref{sec:diff_prog}). Finally, we apply these tools to the Lorenz-63 system to characterize the behavior of parameter gradients in the presence of chaos and discuss the implications for calibrating full-scale climate models (Section~\ref{sec:gradient_unreliability}).

\section{Chaos in Deterministic Dynamical Systems}
\label{sec:chaos}
A continuous-time dynamical system describes the time evolution of a state vector $x(t) \in \mathbb{R}^n$. The governing equations of such systems can take many forms, including those that incorporate explicit time-dependence or stochastic forcing. However, the foundational properties of deterministic chaos are traditionally studied within the specific framework of autonomous deterministic ordinary differential equations (ODEs) of the form $\dot{x} = F[x(t)]$. Because an autonomous system does not explicitly depend on time, its future evolution is entirely dictated by its current state. In this setting, for each initial state $x(0)$, the exact solution $x(t)$ traces a unique trajectory, and the mapping $\phi_t:x(0)\mapsto x(t)$ defines the flow of the system \parencite[Ch.~1.3]{ott_chaos_2002}.

Although the flow $\phi_t$ generated by an autonomous ODE is entirely deterministic, nonlinear systems can produce irregular aperiodic behavior typically referred to as chaos. There are several non-equivalent notions of chaos in the dynamical-systems literature, so it is useful to highlight a few commonly appearing attributes:

\begin{itemize}
    \item \textbf{Sensitive dependence to initial conditions}, where infinitesimal perturbations grow exponentially.
    \item \textbf{Topological Entropy}, quantifying the growth of different orbits,
    \item \textbf{Transitivity with dense orbits}, meaning that trajectories eventually explore the entire accessible phase space
    \item \textbf{Expansivity}, in which distinct trajectories are forced apart under the flow.
\end{itemize}
See \textcite{effah-poku_study_2018} for a detailed discussion. 

A widely used working definition, following \textcite[Sec. 9.3]{strogatz_nonlinear_2018} is:
\begin{quote}
\textbf{Chaos} is \emph{aperiodic long-term behavior} in a \emph{deterministic} system that 
shows \textbf{sensitive dependence on initial conditions}.
\end{quote}

Sensitivity dependence means that trajectories starting arbitrarily close can diverge exponentially fast. A way to measure this is through the largest Lyapunov exponent $\lambda = \lambda_{\max}$, which describes the exponential growth rate of infinitesimal perturbations in generic directions. While an n-dimensional system possesses a full spectrum of $n$ Lyapunov exponents, corresponding to the eigen-directions in phase space and typically including both positive and negative values, the maximal exponent $\lambda_{\max}$ dominates long-term divergence of trajectories. \parencite[Sec.~4.4]{ott_chaos_2002}. A positive $\lambda_{max}$ drives the exponential divergence of the nearby states. This dependence of sensitivity on initial conditions results in deterministic chaos if the system is strictly bound in phase space, forcing exponentially divergent trajectories to continuously fold and mix, rather than diverging to infinity \parencite[Ch.~9.3]{strogatz_nonlinear_2018}.  

For two nearby trajectories with initial separation $\delta_0$, the linearized dynamics gives
\[
\|\delta(t)\| \approx \|\delta_0\|\, e^{\lambda t},
\]

Because the separation increases as $\|\delta(t)\|\sim \|\delta_0\| e^{\lambda t}$, the predictions remain meaningful only while the separation remains below a given tolerance a. The breakdown time, often called the predictability horizon in atmospheric sciences, $t_{\text{horizon}}$ satisfies $\|\delta(t_{\text{horizon}})\| = a$, which gives the following.
\[
t_{\text{horizon}} \approx \frac{1}{\lambda} \ln(\frac{a}{\|\delta_0\|}).
\]
as shown in \textcite[Ch.~9.3]{strogatz_nonlinear_2018}. 

The factor $\frac{1}{\lambda}$ sets the characteristic predictability timescale. The dependence on the initial uncertainty is only logarithmic: improving $\|\delta_0\|$ by many orders of magnitude lengthens the prediction window only by an additive amount of the order $\ln(\|\delta_0\|^{-1})$. This weak sensitivity is the reason why long-term trajectory prediction is impossible in chaotic systems \parencite[Ch.~9.3]{strogatz_nonlinear_2018}. 

\subsection{The Lorenz 63 System}
\label{sec:lorenz63}

\begin{figure}[h]
    \centering
    \includegraphics[width=\textwidth]{thesis/figures/02_lorenz63_chaos_gradients/lorenz_attractor_multipanel_publication.pdf}
    \caption{
        \textbf{The classic Lorenz attractor} generated with parameters 
        $\sigma = 10$ (Prandtl number), $\rho = 28$ (Rayleigh number), 
        and $\beta = 8/3$ (geometric factor). 
        The trajectory was computed from the initial condition $(1,1,1)$ 
        over $100$ time units using a 4th-order Runge–Kutta integrator 
        with timestep $\Delta t = 0.005$. 
    }
    \label{fig:lorenz_attractor}
\end{figure}

In atmospheric physics, a commonly used minimal model exhibiting chaotic behavior is the Lorenz 63 system. It arose from the attempt by Edward Lorenz to understand why even simple models of atmospheric convection could produce irregular, unpredictable motion \parencite{lorenz_deterministic_1963}. Starting from a truncated Fourier expansion of the Rayleigh-Bernard convection equations by \textcite{saltzman_finite_1962}, Lorenz reduced the model to three coupled nonlinear ODEs, as shown in Equation \ref{fig:l63}.

\begin{equation}
\begin{cases}
\dot{x} = \sigma (y - x),\\[4pt]
\dot{y} = x(\rho - z) - y,\\[4pt]
\dot{z} = xy - \beta z,
\end{cases}
\label{fig:l63}
\end{equation}

Physically, the three variables describe the convection of heat. X represents the strength of the convective motion, y is related to the horizontal temperature contrast between ascending and descending air, and z measures the deviation of the vertical temperature profile from a linear gradient \parencite[Ch.~9.1]{alligood_chaos_1996}. The parameters $\sigma$ and $\rho$ correspond to the Prandtl number and a scaled Rayleigh number, while $\beta$ is a geometric factor arising from the convection cell \parencite{lorenz_deterministic_1963}. Commonly chosen parameters are $\sigma = 10,\ \rho = 28,\ \beta = \tfrac{8}{3}$, for which the system produced the classical Lorenz attractor: a bounded, non periodic orbit whose trajectories alternate between two lobes in an apparently irregular fashion, as seen in Figure \ref{fig:lorenz_attractor}. We can compute its Lyapunov exponents numerically, yielding the spectrum $\lambda_1 \approx 0.906 > 0$, $\lambda_2 = 0$, $\lambda_3 \approx -14.572$ \parencite{viswanath_lyapunov_1998}. Time is measured in model time units. The integration variable $t$ in Equation~\ref{fig:l63} is dimensionless, so each numerical step increases the time by $\Delta t = 0.005$, and thus a trajectory of $N$ steps corresponds to a total duration of $T = N \Delta t$ time units.

The Lorenz-63 system provides a controlled setting for studying gradient-based parameter estimation under chaotic dynamics. The goal of parameter calibration here is to recover $\theta = (\rho, \sigma, \beta)$, the set of parameters governing the system, from an observed trajectory $\bar{x}(t)$ generated by true (but unknown) parameters $\theta^* = (10,28,8/3)$. Using only 3 parameters, this serves as a minimal parameter inference problem otherwise similar to full climate models, where the governing equations are known, but the values of physical parameters must be tuned to best reproduce observations.

A differentiable scalar objective $\mathcal{L}(\theta)$ quantifies the mismatch between the model and the reference data. Gradient descent and its variants minimize this objective by iteratively updating the parameters in the direction of steepest descent. This direction is dictated by the negative gradient, which yields the fundamental update rule

\begin{equation}
    \theta_{t+1} = \theta_t - \eta \nabla_\theta \mathcal{L}(\theta_t),
\end{equation}

where $\eta>0$ is the step size or the learning rate \parencite{blondel_elements_2024}.

To apply the gradient-based updates above, we need to compute $\nabla_\theta \mathcal{L}(\theta)$. Since the trajectory $x_\theta(t_k)$ is generated by numerically integrating the Lorenz system with parameters $\theta$, differentiating $\mathcal{L}$ with respect to $\theta$ requires differentiating \emph{through} the integrator. In this work, we compute these derivatives using automatic differentiation within a differentiable-programming framework.

\section{Differentiable Programming}
\label{sec:diff_prog}

Differentiable programming is a way to write numerical models so that their derivatives can be obtained \emph{automatically} and with little overhead, using \emph{automatic differentiation (AD)} instead of hand-written tangent-linear or adjoint code \parencite{baydin_automatic_2015}. In such a setup, once a loss functional is specified, gradients for a very large parameter set can be computed at a computational expense that scales roughly linearly with a single forward model run, up to a small constant \parencite{baydin_automatic_2015,blondel_elements_2024}. This excellent scaling with respect to parameter count explains why differentiable programming powers current deep-learning libraries, where gradient-based methods, like gradient descent and its variants, train massive neural networks. For Earth system models (ESMs), the approach offers comparable advantages as it supports gradient-based tuning of physical parameters, rigorous uncertainty assessment, and hybrid ESMs where machine learning components provide support for hard-to-tune or expensive physics \parencite{gelbrecht_differentiable_2023}.

This parameter-independent scaling stands in stark contrast to the traditional objective calibration methods discussed in Chapter \ref{chap:general-tuning}. For example, while the One-At-a-Time sensitivity analysis scales linearly with the number of parameters as $\mathcal{O}(MN)$, where $M$ is the number of values tested per parameter and $N$ is the total number of parameters, a comprehensive full factorial perturbed parameter ensemble requires a number of model evaluations that scale exponentially as $\mathcal{O}(M^N)$, which quickly becomes computationally infeasible even for a small number of parameters \parencite{eidhammer_extensible_2024}. Computing gradients via finite differences for deterministic inverse problems (e.g., Gauss-Newton algorithms) requires a computational cost that scales between $\mathcal{O}(P)$ and $\mathcal{O}(P^2)$, where $P$ is the number of tuned parameters \parencite{tett_calibrating_2017}. Similarly, the computational cost of surrogate-based tuning and ensemble based data assimilation is contained by sampling frequency. To construct accurate emulators, the required number of model simulations must grow as the number of parameters increases. This can lead to at least ten times as many model runs as parameters \parencite{loeppky_choosing_2009,raoult_parameter_2025,bonnet_tuning_2025}.

Although handwritten adjoint models share the highly efficient scaling properties of automatic differentiation, they impose a massive and nonflexible software engineering burden that can take years to develop manually for complex models \parencite{raoult_parameter_2025}. Instead of requiring developers to manually mirror every forward equation and parameterization change in the adjoint code, differentiable programming systems derive gradients directly from the model low-level code representation, accounting for changes in the forward model \parencite{moses_instead_2020}.

From an implementation standpoint, AD differs from manually computing derivatives, from numerical differentiation, and from symbolic differentiation in that it runs the program in a non-standard way \parencite{baydin_automatic_2015}. The original computation is augmented so that, alongside every numerical operation, the corresponding derivative is carried through. In line with \cite{blondel_elements_2024}, we can view any program as a composition of functions defining a mapping $f : \mathcal{S}_0 \subseteq \mathbb{R}^n \longrightarrow \mathcal{S}_K \subseteq \mathbb{R}^m $: 
$$
f := f_K \circ f_{K-1} \circ \dots \circ f_1.
$$

Each function maps between intermediate states $f_k: \mathcal{S}_{k-1} \to \mathcal{S}_k$, combined by elementary operations. The derivative of $f$ is its Jacobian matrix $J_f(s_0) = \frac{\partial f}{\partial s_0}(s_0) \in \mathbb{R}^{m\times n}$. However, explicitly forming or storing this complete matrix is almost never useful in scientific computing: its size grows as $O(mn)$, while most numerical algorithms depend only on the Jacobian product with a vector, not the individual entries of the matrix. This reduces our interest to two linearized quantities which are computed by the two modes of automatic differentiation.

\begin{itemize}
    \item \textbf{Forward mode}: computes the Jacobian-Vector product (JVP) $J_f(s_0)\,u$ with $u\in\mathcal{S}_0$, giving the directional derivative of f along the input perturbation u.
    \item \textbf{Reverse mode}: computes the Vector-Jacobian product (VJP) $v^\top J_f(s_0)$ with $v\in\mathcal{S}_K$, giving the sensitivity of an output linear function with respect to all inputs
\end{itemize}

\subsection{Forward mode}

Forward mode AD evaluates JVP $J_f(s_0)u$ by propagating tangent perturbations forward through the computational graph. This means that each intermediate state $s_k = f_k(s_{k-1})$ is combined with a tangent $t_k$ that represents the directional derivative of $s_k$ along the perturbed input perturbation $u$. The tangent evolves according to 
\[
 t_k = \partial f_k(s_{k-1}) [t_{k-1}]
\qquad k = 1, \dots, K.
\]
With the initial state $t_0 = u$ the final tangent $t_K$ is the JVP $J_f(s_0)u$. 

This is equivalent to computing with dual numbers $x^{*} =x+\varepsilon, \dot x$ where $\varepsilon^2 =0$: each elementary operation is lifted so that the $\varepsilon$ component tracks the directional derivative at machine precision \parencite{baydin_automatic_2015}.

The forward mode therefore computes a single Jacobian column, making it efficient when the input dimension $n$ is small compared to the output dimension $m$.

\begin{algorithm}[H]
\caption{Forward-mode autodiff algorithm, from \textcite{blondel_elements_2024}}
\begin{algorithmic}[1]
    \Statex \textbf{Functions:} $f := f_K \circ \dots \circ f_1$
    \Statex \textbf{Inputs:} input $\boldsymbol{s}_0 \in \mathcal{S}_0$, input direction $\boldsymbol{u} \in \mathcal{S}_0$
    
    \State Initialize $\boldsymbol{t}_0 := \boldsymbol{v}$
    
    \State \textbf{for} $k := 1, \dots, K$ \textbf{do}
    \State \quad Compute $\boldsymbol{s}_k := f_k(\boldsymbol{s}_{k-1}) \in \mathcal{S}_k$
    \State \quad Compute $\boldsymbol{t}_k := \partial f_k(\boldsymbol{s}_{k-1})[\boldsymbol{t}_{k-1}] \in \mathcal{S}_k$
    
    \Statex \textbf{Outputs:} $f(\boldsymbol{s}_0) := \boldsymbol{s}_K$, $\partial f(\boldsymbol{s}_0)[\boldsymbol{u}] = \boldsymbol{t}_K$
\end{algorithmic}
\end{algorithm}

\subsection{Reverse mode}

Reverse-mode AD evaluates the VJP $v^\top J_f(s_0)$ by propagating adjoint variables for each intermediate variable back through the computation. For each intermediate state $s_k$, the adjoint variable is defined as $r_k = \frac{\partial \mathcal{L}}{\partial s_k}$ for some objective $\mathcal{L}$, which in machine learning is considered the loss function. It represents the sensitivity of the loss function in the intermediate state $s_k$ \parencite{baydin_automatic_2015}. After a forward pass, computing the states $s_1, \dots, s_K$, the adjoint $r_K = v$ is initialized. The adjoint associated with each intermediate variable is then computed via the recursion
\[
r_{k-1} = \partial f_k(s_{k-1})^*[r_k]
\qquad k = K,\dots, 1
\]
where $\partial f_k(s_{k-1})^*$ is the adjoint of $\partial f_k(s_{k-1})$.
The resulting adjoint $r_0$ is the desired VJP $v^\top J_f(s_0)$.

Reverse mode is particularly efficient when the output dimension $m$ is much smaller than the input dimension $n$. In many loss functions and parameter calibration, the output is typically a single scalar loss ($m=1$) that measures the model error, while the inputs are the tunable parameters ($n$). Reverse-mode AD computes the derivative with respect to all $n$ inputs during a single backward pass, so its cost does not increase with parameter count \parencite{baydin_automatic_2018,blondel_elements_2024}. As a result, obtaining the full gradient requires exactly one reverse sweep regardless of the size of the problem, be it calibrating the 3 parameters of Lorenz-63, the dozens to hundreds of physical coefficients in a climate model \parencite{gelbrecht_differentiable_2023}, or the millions of weights in a neural network. In this thesis, the loss is always a scalar trajectory or climatology metric ($m=1$), with $n=3$ for Lorenz-63 and $n=15$ for SpeedyWeather.jl shortwave radiation tuning. Reverse-mode AD thus delivers the complete gradient vector $\nabla_\theta \mathcal{L}(\theta)$ in a single pass, making optimization feasible.

\begin{algorithm}[H]
\caption{Reverse-mode autodiff, from \textcite{blondel_elements_2024}}
\begin{algorithmic}[1] % The [1] turns on line numbering
    \Statex \textbf{Functions:} $f := f_K \circ \dots \circ f_1$,
    \Statex \textbf{Inputs:} input $\boldsymbol{s}_0 \in \mathcal{S}_0$, output direction $\boldsymbol{v} \in \mathcal{S}_K$
    
    \State for $k := 1, \dots, K$ do
    \Comment{Forward pass}
    \State \quad Compute $\boldsymbol{s}_k := f_k(\boldsymbol{s}_{k-1}) \in \mathcal{S}_k$
    
    \State Initialize $\boldsymbol{r}_K := \boldsymbol{v}$.
    
    \State for $k := K, \dots, 1$ do
    \Comment{Backward pass}
    \State \quad Compute $\boldsymbol{r}_{k-1} := \partial f_k(\boldsymbol{s}_{k-1})^*[\boldsymbol{r}_k] \in \mathcal{S}_{k-1}$
    
    \Statex \textbf{Outputs:} $f(\boldsymbol{s}_0) := \boldsymbol{s}_K$, $\partial f(\boldsymbol{s}_0)^*[\boldsymbol{v}] = \boldsymbol{r}_0$
\end{algorithmic}
\end{algorithm}

\subsection{Implementation of AD in Julia: The Enzyme project}

Automatic differentiation becomes significantly more challenging when applied to actual code rather than to an ideal computational chain. In machine learning, for example, AD is typically executed on code that is written in a purely functional array-oriented style. Each operation is an elementary tensor transformation, so that the entire program can be represented as a directed acyclic graph (DAG) \parencite{baydin_automatic_2015}. 

Scientific code rarely has this form. Atmospheric and oceanic models for example, even in a simplified form, rely extensively on in-place updates, time stepping loops, heterogeneous mutable data structures, and branching logic. These features are used primarily for memory efficiency as, for example, in-place updates avoid the allocation of new arrays. However, they introduce complex memory dependencies and mutations that AD systems must differentiate alongside mathematical expressions, not just high-level equations.

Enzyme \parencite{moses_instead_2020} is the AD engine used in this work. Unlike most AD frameworks, Enzyme does not operate at the level of Julia source code. Instead, it acts on the program after the Julia compiler has produced the LLVM intermediate representation (IR). LLVM IR is a low-level assembly-like language that encodes the control-flow graph (CFG), data dependencies, memory operations, and the effects of mutable updates \parencite{noauthor_llvm_nodate}. Because LLVM has already performed an optimization step on the IR before passing it on to Enzyme, it is significantly simpler than the original code. By working at this level, Enzyme is able to differentiate Julia code without requiring the user to rewrite the model into a functional style.

\subsection{Differentiation at the compiler IR level}

Julia compiles functions through a multi-stage pipeline to progressively lower high level abstractions to machine code. The process begins with parsing, where the source code is converted into an abstract syntax tree (AST). This AST is then lowered to a normalized lowered form in which control flow constructs, like blocks and jumps, are explicit. This lowered representation is then passed to a typed inference engine, which performs abstract interpretation to produce a typed intermediate representation (IR) following a static single-assignment (SSA) style. This typed IR is used to generate LLVM IR, a low-level, language independent representation, and then undergoes an optimization step. At this stage, the program is represented using explicit memory operations, typed primitives, and control-flow blocks \cite{julia_developers_julia_nodate}.

Enzyme analyzes the IR translation of the user code, constructs a derivative version of the functions, and gives the new IR to the compiler for further optimization. This derivative version depends on which differentiation mode is chosen. In a sense, AD becomes a compiler pass. The differentiation occurs before machine code generation, ensuring that the resulting derivative undergoes the same optimization as the primary function \parencite{moses_instead_2020}. This architecture gives Enzyme two key advantages. First, it can differentiate through arbitrary control flows. Conditionals and loops appear explicitly in LLVM IR. Enzyme therefore has access to the full control flow graph. Differentiation a branch corresponds to differentiating the branch condition and the selected basic block, while loops are handled by reversing their control flow structure in reverse mode. Secondly, because IR includes explicit memory operations, Enzyme can track adjoints of loads and stores, thereby differentiating through in-place mutations.

\subsection{The Adjoint Generation Process}

Enzyme applies a sequence of compiler analyzes and transformations to generate a correct and efficient adjoint for a given LLVM function \parencite{moses_instead_2020}. It begins by inferring types and activity information on the LLVM IR, then inserts shadow memory and tape data structures, and ultimately constructs and optimizes a reverse pass that can be invoked as any other compiled function \parencite{moses_instead_2020}.

\paragraph{Type Analysis.}
LLVM pointer and memory operations do not always preserve the high-level type of the underlying data. Code often uses generic pointers and operations, so the IR type alone does not reveal whether a location stores an integer, a floating point value, or a pointer \parencite{moses_instead_2020}. To generate semantically correct adjoints, Enzyme implements a custom type analysis that builds a type tree for every SSA value and memory location, propagating information from the allocation sites and uses \parencite{moses_instead_2020}. This allows Enzyme to distinguish pointer arithmetic from floating-point computation, avoid taking derivatives with respect to non-differentiable data such as integer loop indices, and align shadow values with the correct primal types \parencite{moses_instead_2020,noauthor_llvm_nodate}.

\paragraph{Activity Analysis.}

Building on the recovered type information, Enzyme then identifies which values and instructions are \emph{active}, meaning those that can affect the specified objective and thus need derivatives \parencite{moses_instead_2020,peng_sound_2025}. This activity analysis operates on LLVM control-flow graphs and labels SSA values, as well as the memory locations they occupy, as either active or inactive, while properly handling aliasing \parencite{peng_sound_2025}. Only active computations are included in the adjoint, avoiding unnecessary differentiation of purely auxiliary control-flow or bookkeeping code, which is crucial for both correctness and efficiency.

\paragraph{Shadow Memory and Tape Synthesis.}

For each active value produced during the forward pass, Enzyme allocates a corresponding \emph{shadow} slot to hold its derivative, for active data structures it creates parallel shadow structures that replicate the layout of the original (primal) memory \parencite{moses_instead_2020,moses_reverse-mode_2021}. Moreover, Enzyme must guaranty that all values required for an optimal reverse pass remain accessible, even if they would otherwise be overwritten during the forward computation. To achieve this, it generates a \emph{tape} that records selected intermediate quantities, balancing recomputation cost against storage use \parencite{moses_instead_2020}. This tape is implemented directly in LLVM IR as additional arrays and temporaries, making its behavior fully explicit to subsequent optimization passes \parencite{moses_instead_2020}.

\paragraph{Reverse Pass Generation.}

Once type and activity information, shadow memory, and tape layout are set, Enzyme synthesizes the reverse pass by transversing each basic block in reverse order and emitting the corresponding adjoint instructions \parencite{moses_instead_2020,moses_supercharging_nodate}. For every primal instruction, Enzyme applies a differentiation rule that updates the appropriate shadow values, loads any cached primal operands from the tape if necessary, and respects the original control flow by creating mirrored reverse blocks and branches \parencite{moses_instead_2020}. The resulting gradient function combines the original forward computation, shadow initializations, and reverse sweep into a single LLVM function, or into paired forward and reverse functions when required and aliasing constraints \parencite{moses_enzyme_2020}.

\paragraph{Optimization and Lowering}

Because adjoint synthesis is implemented as a compiler pass on LLVM IR, the generated gradient code can then be optimized by standard LLVM pipelines in the same way as the primal function \parencite{moses_instead_2020,moses_enzyme_2020}. Classical optimization techniques apply equally to reverse pass and tape accesses \parencite{moses_enzyme_2020}. Finally, both primal and adjoint IR are lowered together to machine code for the target architecture, so that, on the Julia side, Enzyme-based differentiation exposes gradients as ordinary callable functions integrated into the usual compilation flow \parencite{moses_reverse-mode_2021}.

\section{Parameter estimation in Lorenz-63 via AD}
\label{sec:param_estimation}

We use Lorenz-63 as a minimal inverse problem to answer the question if reverse-mode AD can recover model parameters from trajectory data in a chaotic system and if so over which forecast horizon. The reference trajectory is generated from the initial condition $\bar{x}(0)=(1,1,1)$ with the true parameters $\theta^\star=(\sigma^\star,\rho^\star,\beta^\star)=(10,28,8/3)$ on $T=100\,\mathrm{t.u.}$ using $\Delta t=0.005$, corresponding to $M=20{,}000$ time steps. The experiments were carried out in three stages: first as a short-horizon proof of concept that recovers only $\rho$, then as simultaneous recovery of all three parameters, and finally as a systematic study of how the useful gradient signal breaks down as the trajectory window approaches longer timescales. The central result is that reverse-mode AD itself is not the obstacle. The obstacle is the horizon over which the trajectory loss is defined, and thus the gradient is defined.

A naive formulation would keep the initial condition fixed and minimize a full-trajectory mismatch,
\begin{equation}
    \mathcal{L}_{\mathrm{full}}(\theta) = \frac{1}{M}\sum_{k=0}^{M-1}\left\|x_\theta(t_k)-\bar{x}(t_k)\right\|^2.
\end{equation}

In a chaotic system, this is problematic as even when $\theta$ is close to $\theta^\star$, the state error grows with the divergence of the trajectories. So after about one Lyapunov time the simulated orbit decorrelates from the reference. Beyond that point, the loss is dominated by the attractor-scale separation between two unrelated trajectories, rather than by the original parameter error.

Figure~\ref{fig:lorenz_gradient_uncertainty} shows the $\ell_2$ norm of the trajectory-loss gradient calculated at the perturbed parameter vector $(\sigma,\rho,\beta)=(10,25,8/3)$ and averaged over 30 randomly chosen window starts. As the window length increases from $0.23\,\tau_L$ to $20.4\,\tau_L$, the mean gradient norm increases by around five orders of magnitude, and the variance across also widens strongly. Long-horizon differentiation therefore does not yield a stable descent direction, it yields a high-variance object whose value depends sensitively on the chosen segment of the chaotic orbit.

\begin{figure}[h]
    \centering
    \includegraphics[width=1\linewidth]{thesis//figures//02_lorenz63_chaos_gradients/lorenz_gradient_uncertainty.pdf}
    \caption{\textbf{Exponential growth of gradient norms with window length in Lorenz-63.} Mean $\pm 1\sigma$ of $\log_{10}\lVert\nabla_\theta\mathcal{L}_j\rVert_2$ across 30 random window starts, evaluated at the perturbed parameter vector $(\sigma,\rho,\beta)=(10,25,8/3)$ against a reference trajectory at the true parameters $(\sigma^*,\rho^*,\beta^*)=(10,28,8/3)$. The dotted reference line shows the expected growth rate $e^{\lambda t}$ with the true Lyapunov exponent $\lambda=0.906$. The shaded ribbon (one standard deviation across windows) widens with $T_w$: for $T_w\lesssim\tau_L$ gradient estimates are consistent across windows, while beyond $\tau_L\approx1.10$ (vertical line) the spread exceeds one order of magnitude, making individual gradient directions unreliable for optimization.
}
    \label{fig:lorenz_gradient_uncertainty}
\end{figure}

The solution is to reformulate the parameter estimation into as many short forecasting tasks. As illustrated in Figure~\ref{fig:p12_windows_overview}, the reference trajectory is split into fixed lengths of (possibly) overlapping windows. For each window $j$, the model is reinitialized with the observed start state $\bar{x}^{(j)}(t_0)$ and integrated only over that short segment, a standard strategy known as \emph{teacher forcing}. The per-window loss is

\begin{equation}
    \mathcal{L}_j(\theta) = \frac{1}{M_w}\sum_{k=0}^{M_w-1}\left\|x_\theta^{(j)}(t_k)-\bar{x}^{(j)}(t_k)\right\|^2,
    \label{eq:mse_loss}
\end{equation}
where $M_w$ is the window length in time steps. Restarting from the observed state prevents exponential drift from accumulating across the full trajectory and keeps the loss informative about the parameters over the duration of the window.

\begin{figure}[h]
    \centering
    \includegraphics[width=1\linewidth]{thesis//figures//02_lorenz63_chaos_gradients/p12_windows_overview.png}
    \caption{\textbf{Windowed trajectory sampling for stable gradient-based parameter recovery in the chaotic Lorenz-63 system.}
    (a) Lorenz-63 attractor with trajectory windowing. The reference trajectory (gray, background) is partitioned into overlapping windows of 300 time steps (colored paths: W1--W5 highlighted). Each window provides a short-horizon training example for gradient-based parameter recovery via Enzyme reverse-mode AD. 
    (b) Window overlap timeline (window=$300$ steps, stride=$150$ steps, $50\%$ overlap). Horizontal bars show temporal extent of first $9$ windows; vertical dotted lines mark stride boundaries. Overlap ensures sampling of the chaotic attractor while limiting backpropagation horizons below the Lyapunov time. This mini-batch structure enables stable gradient averaging across $B=4$ windows per update in $`modular\_train!`$.}
    \label{fig:p12_windows_overview}
\end{figure}

At each update, a batch $\mathcal{B}$ of $B$ windows is sampled and the parameter update is based on the average gradient

\begin{equation}
    g(\theta) = \frac{1}{B}\sum_{j\in\mathcal{B}} \nabla_\theta \mathcal{L}_j(\theta).
    \label{eq:minibatch_gradient}
\end{equation}

This averaging step is essential. A single window only samples one local segment of the attractor and, therefore, yields a noisy trajectory-local sensitivity. Averaging over several windows converts many of such local sensitivities into a useful update direction.

The entire procedure is implemented in the \texttt{modular\_train!\footnote{By Julia convention, functions ending in \texttt{!} mutate at least one of their arguments in place rather than returning a new value.}} routine of the LorenzParameterEstimation.jl package (\textcite{niklas_viebig_nviebiglorenzparameterestimation_2026}). At each training step, the algorithm (i) samples a batch of window start indices, (ii) extracts the corresponding observed initial conditions, (iii) integrates a short RK4 forecast in each window, (iv) differentiates the window loss with Enzyme in reverse mode, (v) averages the gradients across the batch, and (vi) applies a first-order optimizer to the active parameters.

For one window, the reverse-mode differentiation call takes the form

\begin{lstlisting}[
caption={Reverse-mode differentiation call in Enzyme.jl\protect\parencite{moses_instead_2020}},
label={lst:enzyme_call},
mathescape=true
]
Enzyme.autodiff(
    Enzyme.Reverse,
    L,
    Active($\sigma$), Active($\rho$), Active($\beta$),
    Const($\bar{x}_0$),
    Const($\bar{x}_{1:M_w}$),
    Const($M_w$),
    Const($\Delta t$)
)
\end{lstlisting}

where \texttt{Active} indicates the differentiable parameters and \texttt{Const} indicates the observed start state, reference window, and numerical settings. The result is a tuple $(L,(g_\sigma,g_\rho,g_\beta))$, so a reverse sweep returns both the window loss and all parameter gradients at a cost of roughly one additional forward integration.

\paragraph{Optimizer and gradient clipping.}

Adam \parencite{blondel_elements_2024} is used as the default optimizer, with learning rate $\eta=0.01$ and momentum coefficients $\beta_1=0.9$, $\beta_2=0.999$. This is helpful because $\sigma$, $\rho$, and $\beta$ induce gradients of different natural scales. Before the update, the batch-averaged gradient is clipped in the $\ell_2$ norm to prevent rare large windows from destabilizing optimization. Separate optimizer comparisons show that, once the window is short enough, SGD, AdaGrad, and Adam all converge reliably; the dominant design choice is the window length, not the optimizer family. With 100-step windows, the mean final parameter-error norm over five seeds lies between $8.9\times10^{-4}$ and $2.3\times10^{-3}$ across these optimizers (Appendix~\ref{app:additional-figures}, Figure~\ref{fig:optimizer_comparison_publication}).

The experiments were intentionally staged. A first proof-of-concept estimated only $\rho$, recovering it from an initial estimate of $\rho_0=17$ to approximately $\rho^\star=28$ while keeping $\sigma$ and $\beta$ fixed. This verified that Enzyme gradients through the RK4 loop have the correct sign and scale. The fully coupled experiment then allowed all three parameters to vary simultaneously.

\begin{figure}[h]
    \centering
    \includegraphics[width=0.8\linewidth]{thesis//figures//02_lorenz63_chaos_gradients/lorenz_training_convergence.png}
    \caption{
    \textbf{Trajectory-based recovery of all Lorenz-63 parameters with short windows.} Panels (a)--(c) compare trajectories integrated from the initial parameter guess $(\sigma,\rho,\beta)=(6.0,20.0,1.0)$, the ground-truth parameters $(10.0,28.0,8/3)$, and the recovered parameters after 60 epochs. Panels (d)--(f) show the corresponding training diagnostics: the windowed RMSE, the parameter trajectories, and the absolute parameter errors. Gradients are computed over overlapping 100-step windows (stride 50, batch size 8) using reverse-mode Enzyme differentiation. In this short-window regime the optimization converges smoothly despite the underlying chaotic dynamics.}
    \label{fig:lorenz_training_convergence}
\end{figure}

Figure~\ref{fig:lorenz_training_convergence} shows a recovery run of all-parameters. Starting from the biased initial vector $(6,20,1)$, the short-window formulation drives the parameters back toward $(10,28,8/3)$ and steadily reduces the loss. The key point is that the trajectory-matching problem is solvable as long as it is posed as many short forecasts rather than as one long chaotic roll-out.


Once the short-window formulation is in place, the remaining question is how short the window must be. An analysis is given in Appendix~\ref{app:experiments}, where all three parameters are recovered simultaneously over eight window lengths that range from $0.1\,\tau_L$ to $20\,\tau_L$ and 100 random seeds (Figure~\ref{fig:all_params_window_experiment}). The transition is sharp. For windows up to about $1\,\tau_L$, the mean final error remains very small and is tightly concentrated. At $1\,\tau_L$ the mean errors are $|\hat{\sigma}-\sigma^\star|\approx 0.009$, $|\hat{\rho}-\rho^\star|\approx 0.004$, and $|\hat{\beta}-\beta^\star|\approx 0.002$. Beyond roughly one Lyapunov time the final errors rise rapidly and the spread across seeds widens strongly. By $2\,\tau_L$ and above, many runs no longer recover the correct parameter vector.

The lower bound is equally interesting, as even the shortest window tested, the $0.1\,\tau_L \approx 18$--$22$ time steps, already carries useful parameter information because $\sigma$, $\rho$, and $\beta$ enter the Lorenz right-hand side directly. This indicates that for Lorenz-63 the binding constraint is the upper limit imposed by chaos, not the lower limit required for a signal to appear. This will not be true in complex climate models, where some parameters only affect the loss after much longer physical adjustment times.

The loss-landscape analysis explains this transition geometrically. We compute $\mathcal{L}_j(\theta)$ on a $21\times 21$ grid over each pair of parameters while keeping the third parameter fixed at its true value, using three random windows per point on the grid.

\begin{figure}[h!]
    \centering
    \includegraphics[width=1\linewidth]{thesis//figures//02_lorenz63_chaos_gradients/landscape_2d_selected.pdf}
    \caption{
    \textbf{Two-dimensional loss landscapes for Lorenz-63 parameter estimation.}
    Each panel shows $\log_{10}(\mathrm{RMSE})$ on a regular grid over two parameters while the third is held fixed at its true value. \textit{Rows}: parameter pairs $(\sigma,\rho)$, $(\beta,\rho)$, and $(\sigma,\beta)$. \textit{Columns}: trajectory window lengths of $0.1\,\tau_L$, $1.0\,\tau_L$, and $10.0\,\tau_L$, where $\tau_L \approx 1.10\,\mathrm{t.u.}$. The \textbf{red star}~($\bigstar$) marks the true parameter value $(\sigma^{\!*}\!=10,\,\rho^{\!*}\!=28,\,\beta^{\!*}\!=8/3)$; a \textbf{red diamond} marks the initial guess.}
    \label{fig:landscape_2d_selected}
\end{figure}

Figure~\ref{fig:landscape_2d_selected} shows three representative horizons in all three parameter planes. The complete set is shown in Figure~\ref{fig:loss_landscapes_all}. At $0.1\,\tau_L$ the landscape is smooth and dominated by a broad basin around the true minimum. Around $1\,\tau_L$ this basin persists, but ridges begin to form, especially in the $\rho$ direction, so gradient descent still works, but becomes less regular. At $10\,\tau_L$ the surface fragments into disconnected low-loss regions with many local minima of comparable depth. In that regime, a single-window gradient mostly measures orbit separation rather than the global parameter dependence of the attractor.

These experiments lead to a clear conclusion. Reverse-mode AD can recover Lorenz-63 parameters, but only when the calibration problem is posed in a short, weather-like regime: many restarted forecasts, each shorter than roughly one Lyapunov time, combined through batch averaging. Once the horizon is extended, the gradients become dominated by chaotic orbit separation and the full-trajectory loss ceases to be a reliable optimization target. Preliminary attempts to replace the trajectory mismatch with PDF-based losses were also made, but within this project they did not yet provide a comparably robust alternative. The main lesson for the remainder of this thesis is therefore that fast processes can still be calibrated from short-horizon trajectory information, whereas slower climate feedbacks require losses based on statistics, distributions, or climatological means rather than on pointwise agreement along individual chaotic trajectories.
